\chapter{Introducere}\label{ch:intro}
\pagestyle{fancy}

\section{Contextul temei}

\textit{"Alert fatigue"} este fenomenul prin care un operator
expus la un volum mare de alerte, dintre care majoritatea
se dovedesc nerelevante, își pierde treptat capacitatea de a 
reacționa corect la fiecare~\cite{Ndichu2026}. Fenomenul
e documentat în securitatea cibernetică, unde analiștii
primesc mai multe alerte decât pot investiga, iar cele mai multe
sunt fals pozitive~\cite{Hassan2019}. Consecințele nu sunt
doar teoretice: atacul asupra companiei Target din 2013, 
prin care s-au furat 40 de milioane de înregistrări
de card, s-a produs parțial pentru că personalul nu a reacționat la timp la niște alerte
primite, clasificându-le greșit drept fals pozitive obișnuite~\cite{Hassan2019}.
Organizațiile
primesc în medie 17.000 de alerte pe săptămână, dintre care
peste 51\% sunt fals pozitive, iar doar 4\% ajung să fie investigate 
corespunzator~\cite{Hassan2019}.

Efectul \textit{"cry-wolf"} descrie scăderea ratei de răspuns la
alerte pe măsură ce crește proporția alertelor false primite~\cite{bliss1995reversal}. 
Un atacator poate genera
în mod repetat alerte false pentru a obosi operatorul, până în punctul în care
o alertă reală este tratată cu aceași superficialitate ca precedentele
false. 

În monitorizare video, sistemul de detecție generează
o alertă însoțită de un clip scurt, iar un operator uman
decide dacă este o intruziune reală~\cite{Lohani2022}.
Articolele dedicate detecției intruziunilor prin supraveghere video
menționează printre alte provocări și alarmele false produse
de variațiile de iluminare, de vegetație sau de animale~\cite{Lohani2022}.
Situația se aseamănă cu gestionarea alertelor într-un centru de
operațiuni de securitate (SOC), dar aplicată fluxurilor video, nu evenimentelor
de rețea.

Când vine vorba de monitorizare video, un posibil atacator poate folosi
în mod tactic vulnerabilitățile sistemului de detecție, declanșând în mod
intenționat alerte repetate pe camerele de supraveghere. Această 
supraîncarcare duce la oboseala operațională a agentului de monitorizare, 
apare fenomenul de \textit{"alert fatigue"}, determinându-l să intre 
într-un mod de funcționare de tip
\textit{"survival"}: verifică rapid, superficial sau ignoră direct, pentru a face față
ritmului impus. Riscul direct al acestui comportament e ca alerte, care
marchează incidente reale, ajung tratate la fel de superficial ca cele false.
Pentru a combate această
tehnică de evaziune, este esențială implementarea unui filtru inteligent
de analiză video, capabil să valideze alertele și să transmită operatorului
doar incidentele ce merită analizate, detaliat în capitolul următor. 

\section{Motivația personală}

Interesul meu pentru această temă provine din activitatea pe care o 
desfășor ca operator în echipa de audit 24/7 a firmai BigBrother International.
Printre sarcinile echipei se numără analiza alertelor declanșate
de sistemul de supraveghere. Fluxul actual de lucru: sistemul de detecție
generează o alertă pe baza unui eveniment observat în fluxul video, iar
alerta ajunge la operator, care trebuie să decidă rapid dacă situația
necesită intervenție (contactare client, echipa de securitate, 
autoritățiile, în funcție de situație si procedură).
Pe langă asta, încercam să înțelegem de ce a fost declanșată fiecare 
alertă, pentru a comunica propuneri de îmbunătățire a sistemului.

Pe durata activității mele am observat o problemă cu efect direct
asupra calității monitorizării. În timpul zilei, sistemul actual
generează un număr foarte mare de alerte, deoarece semnalează
mișcare într-o zonă marcată, indiferent de cauză. Așadar, sistemul nu face
diferența între o alertă reală si una falsă, produsă de exemplu, de 
un animal sau o persoana care doar trece prin cadru.
Cum riscul unui furt sau al unui act de
sabotaj este perceput ca fiind mai redus ziua, iar aproape toate alertele
sunt false, echipa a ajuns treptat să le trateze cu mai puțină atenție. 


Consecințele au fost vizibile în două situații reale.
În prima situație, 
un client a dorit să testeze sistemul. A declanșat în mod
intenționat opt alerte de la camere diferite montate pe locația
sa, apoi s-a apropiat de un cablu, simulând ca vrea să îl taie
cu un fierăstrău. Având de procesat un volum mare de alerte 
și bazându-se pe ideea că în plină zi nu se poate întâmpla nimic
grav, operatorul de tură a închis alertele fără a le vizualiza.
În de a doua situație, un coleg monitoriza un depozit de marfă. 
Odată cu începerea programului de lucru, numărul alertelor a
crescut, iar un furt de paleți comis chiar de către angajații 
locației a trecut neobservat.

Aceste situații nu indică lipsă de profesionalism, ci un efect
previzibil al volumului de alerte asupra atenției umane.
Ele ilustrează riscul descris anterior. Observația din practică
este motivația pentru tema lucrării: un filtru care verifică automat
alertele și transmite operatorului doar incidentele care merită analizate.

\section{Delimitarea domeniului}

Lucrarea se situează în domeniul procesării video asistate
de inteligență artificială, aplicat sistemelor de supraveghere. Sistemul de detecție
existent semnalează evenimente precum intrarea într-o zonă marcată. 
Lucrarea propune un pas suplimentar: un clasificator care primește clipul
video asociat unei alerte deja declanșate și estimează dacă este reală sau falsă, 
înainte ca alerta să ajungă la operator. Cazurile incerte sunt transmise operatorului, 
pentru că un filtru care închide automat o alertă reală ar repeta greșala operatorului obosit.
În această lucrare, o alertă falsă (fals pozitiv) este o
alertă ce nu trebuia trimisă spre analiza operatorului, ea
fiind categorizată drept zgomot.

Soluția nu înlocuiește sistemul de detecție și nu elimină operatorul uman.
Este urmărită reducerea numărului de alerte false pe care operatorul
trebuie să le analizeze. Nu fac obiectul lucrării dezvoltarea unui nou detector, 
decizia de intervenție sau contactarea autorităților. Datele de antrenare
și testare provin din istoricul de alerte și din înregistrările video ale firmei.


\section{Structura lucrării}

Lucrarea este structurată astfel:
\begin{itemize}
    \item Capitolul 2 formulează exact tema și obiectivele proiectului
    \item Capitolul 3 trece în revistă stadiul actual al domeniului
    \item Capitolul 4 fundamentează teoretic soluția propusă
    \item Capitolul 5 descrie proiectarea și implementarea aplicației
    \item Capitolul 6 reprezintă testarea și validarea rezultatelor
    \item Capitolul 7 oferă un manual de instalare și utilizare
    \item Capitolul 8 trage concluziile și conturează direcții viitoare de dezvoltare
\end{itemize}


